\chapter{Scenari Quantizzati e Dashboard
SPC-DEQ}\label{scenari-quantizzati-e-dashboard-spc-deq}

\begin{quote}
\textbf{Argomento portante:} la DEQ$^2$ è una teoria \emph{operativa}
perché ammette una traduzione diretta in \textbf{Statistical Process
Control adattato ai sistemi sociali}. La Dashboard SPC-DEQ è la
proiezione monitorabile di \(\boldsymbol{p}(t) \in \mathbb{V}^{(4)}\) in
\emph{sei pannelli bivariati} (uno per ciascuna relazione R1-R6) più un
\emph{pannello aggregato} della distanza scalare
\(d(\boldsymbol{p}(t), \mathcal{S})\). Le soglie operative non sono
regole arbitrarie ma derivano dalla geometria della superficie critica
\(\mathcal{S}\) e dalla varianza naturale degli indicatori. I tre attori
principali (USA, Cina, Brasile) sono mappati in scenari quantizzati
2026-2031, con uno scenario aggregato e due contro-fattuali brasiliani
che dipendono dall'esito elettorale del 4 ottobre 2026.
\end{quote}



\section{Perché SPC}\label{perchuxe9-spc}

\Hypoth{} \textbf{Tradizione tecnica.} Lo \emph{Statistical Process
Control} (Walter Shewhart, 1924; Edward Deming, 1950s; Six Sigma 1980s)
è la disciplina che monitora processi industriali tramite \emph{carte di
controllo} con limiti \(UCL\) (upper) e \(LCL\) (lower) calibrati sulla
varianza naturale del processo. Le \emph{regole di Western Electric}
(1956) identificano pattern di derive sistematiche prima che il processo
esca dai limiti.

\Hypoth{} \textbf{Adattamento alla DEQ$^2$.} Il sistema sociale è
formalmente analogo a un processo industriale con due differenze:

{\def\LTcaptype{none} % do not increment counter
{\small\begin{longtable}[]{@{}
  >{\raggedright\arraybackslash}p{(\linewidth - 2\tabcolsep) * \real{0.5000}}
  >{\raggedright\arraybackslash}p{(\linewidth - 2\tabcolsep) * \real{0.5000}}@{}}
\toprule\noalign{}
\begin{minipage}[b]{\linewidth}\raggedright
Differenza SPC industriale $\to$ SPC-DEQ$^2$
\end{minipage} & \begin{minipage}[b]{\linewidth}\raggedright
Adattamento richiesto
\end{minipage} \\
\midrule\noalign{}
\endhead
\bottomrule\noalign{}
\endlastfoot
Variabili continue univariate $\to$ variabili strutturali
\(\mathbb{V}^{(4)}\) multivariate & Sei pannelli bivariati (R1-R6) +
pannello aggregato \(d(\boldsymbol{p}, \mathcal{S})\) \\
Limiti \(UCL/LCL\) simmetrici $\to$ superficie critica \(\mathcal{S}\)
asimmetrica & \(\mathcal{S}\) ricavata analiticamente dal Cap. 5.3, non
dalla varianza \\
Causa speciale = fuori-limite $\to$ transizione di fase =
\(T_s^{(2),\text{sys}} < 1\) persistente & Soglia di persistenza
\(\Delta t > \Delta t^*\) (Cap. 3.5.3) \\
Capability index \(C_p = (USL-LSL)/6\sigma\) $\to$ indice DEQ$^2$
\(C_{\text{DEQ}}^2 := d(\boldsymbol{p}, \mathcal{S}) / \sigma_d\) &
Quantità scalare confrontabile fra sistemi \\
\end{longtable}}
}

\Proved{} \textbf{Conseguenza operativa.} La Dashboard SPC-DEQ è
\emph{azione politica formalizzata}: ogni leva del Cap. 5.5 (gradi di
libertà operativi) corrisponde a un intervento mirato su uno o più
pannelli. Non è strumento di previsione passiva --- è \emph{interfaccia
di pilotaggio} del punto \(\boldsymbol{p}(t)\) in
\(\mathbb{V}^{(4)}_{\text{ammissibile}}\).



\section{I Sei Pannelli Bivariati}\label{i-sei-pannelli-bivariati}

\Hypoth{} \textbf{Costruzione.} Per ciascuna relazione \(R_i\) del Cap.
5.1, costruiamo un \emph{pannello bidimensionale} sulle due coordinate
accoppiate, in cui:

\begin{itemize}
\tightlist
\item
  l'asse \(x\) è la prima variabile della coppia;
\item
  l'asse \(y\) è la seconda;
\item
  una \emph{iso-curva} \(\mathcal{S}_i\) rappresenta la sezione di
  \(\mathcal{S}\) a \((s_{iii}, z, \delta, \sigma)\) esogeni fissi e
  altre due variabili al loro valore mediano sistemico;
\item
  il punto attuale \(\boldsymbol{p}(t)\) è proiettato e tracciato come
  \emph{trail temporale} (ultimi 12-24 mesi);
\item
  una \emph{zona di near-miss} (banda colorata tra \(\mathcal{S}_i\) e
  \(0{,}9 \cdot \mathcal{S}_i\) in distanza euclidea normalizzata)
  segnala l'avvicinamento alla criticità.
\end{itemize}

\Proved{} \textbf{I sei pannelli con interpretazione operativa.}

{\def\LTcaptype{none} % do not increment counter
{\small\begin{longtable}[]{@{}
  >{\raggedright\arraybackslash}p{(\linewidth - 8\tabcolsep) * \real{0.2000}}
  >{\raggedright\arraybackslash}p{(\linewidth - 8\tabcolsep) * \real{0.2000}}
  >{\raggedright\arraybackslash}p{(\linewidth - 8\tabcolsep) * \real{0.2000}}
  >{\raggedright\arraybackslash}p{(\linewidth - 8\tabcolsep) * \real{0.2000}}
  >{\raggedright\arraybackslash}p{(\linewidth - 8\tabcolsep) * \real{0.2000}}@{}}
\toprule\noalign{}
\begin{minipage}[b]{\linewidth}\raggedright
Pannello
\end{minipage} & \begin{minipage}[b]{\linewidth}\raggedright
Coppia
\end{minipage} & \begin{minipage}[b]{\linewidth}\raggedright
Iso-curva
\end{minipage} & \begin{minipage}[b]{\linewidth}\raggedright
Direzione di sicurezza
\end{minipage} & \begin{minipage}[b]{\linewidth}\raggedright
Allarme tipico
\end{minipage} \\
\midrule\noalign{}
\endhead
\bottomrule\noalign{}
\endlastfoot
\textbf{P1} & \(\Phi\)-\(A\) &
\(A = (\Phi^{-1} - 1)/\langle A\rangle_{\text{ref}}\) & NE (alto
\(\Phi\), alto \(A\)) & \(\boldsymbol{p}\) scende verso SW (decoerenza
fragilizzante) \\
\textbf{P2} & \(\Phi\)-\(\Omega\) & \(\Phi = \Omega/\Omega_c\) con
\(\Omega_c\) critica & NW (alto \(\Phi\), basso \(\Omega\)) &
\(\boldsymbol{p}\) migra verso SE (decoerenza con dissipazione) \\
\textbf{P3} & \(\Phi\)-\(\varepsilon_{\text{dis}}\) &
\(\Phi = 1 - \varepsilon_{\text{dis}}/\varepsilon_c\) & NW (alto
\(\Phi\), basso \(\varepsilon\)) & \(\boldsymbol{p}\) verso SE
(decoerenza con disimpegno) \\
\textbf{P4} & \(A\)-\(\Omega\) & \(A = -\Omega \cdot k_4\) (\(k_4 > 0\))
& NW (alto \(A\), basso \(\Omega\)) & \(\boldsymbol{p}\) verso SE
(fragilità con dissipazione --- pattern \emph{cinese}) \\
\textbf{P5} & \(A\)-\(\varepsilon_{\text{dis}}\) &
\(A = (1 - \varepsilon)/k_5\) & NW (alto \(A\), basso \(\varepsilon\)) &
\(\boldsymbol{p}\) verso SE (fragilità con disimpegno --- pattern
\emph{USA}) \\
\textbf{P6} & \(\Omega\)-\(\varepsilon_{\text{dis}}\) &
\(\varepsilon_{\text{dis}} = k_6 \cdot \Omega\) & SW (basso \(\Omega\),
basso \(\varepsilon\)) & \(\boldsymbol{p}\) verso NE (dissipazione che
alimenta disimpegno --- pattern \emph{involutivo}) \\
\end{longtable}}
}

\Hypoth{} \textbf{Pannello aggregato \(P_0\).} La distanza scalare
\(d(\boldsymbol{p}(t), \mathcal{S})\) è tracciata su un'unica carta di
controllo con:

\begin{itemize}
\tightlist
\item
  \emph{zona verde}: \(d > 0{,}30\) (margine di sicurezza ampio);
\item
  \emph{zona gialla}: \(0 < d \leq 0{,}30\) (vigilanza);
\item
  \emph{zona rossa}: \(-0{,}20 < d \leq 0\) (sotto soglia ma
  reversibile);
\item
  \emph{zona nera}: \(d \leq -0{,}20\) (transizione di fase consolidata,
  \(\Delta t > \Delta t^*\) richiesto).
\end{itemize}



\section{Soglie Operative e Segnali di Allarme
Precoce}\label{soglie-operative-e-segnali-di-allarme-precoce}

\Hypoth{} \textbf{Adattamento delle regole di Western Electric a
\(\mathbb{V}^{(4)}\).} La logica WE classica si traduce in cinque
pattern DEQ$^2$:

{\def\LTcaptype{none} % do not increment counter
{\small\begin{longtable}[]{@{}
  >{\raggedright\arraybackslash}p{(\linewidth - 4\tabcolsep) * \real{0.3333}}
  >{\raggedright\arraybackslash}p{(\linewidth - 4\tabcolsep) * \real{0.3333}}
  >{\raggedright\arraybackslash}p{(\linewidth - 4\tabcolsep) * \real{0.3333}}@{}}
\toprule\noalign{}
\begin{minipage}[b]{\linewidth}\raggedright
Pattern WE-DEQ$^2$
\end{minipage} & \begin{minipage}[b]{\linewidth}\raggedright
Definizione
\end{minipage} & \begin{minipage}[b]{\linewidth}\raggedright
Interpretazione
\end{minipage} \\
\midrule\noalign{}
\endhead
\bottomrule\noalign{}
\endlastfoot
\textbf{WE-1} & Un punto in zona rossa o nera & Allarme immediato ---
entrata in \(\mathcal{R}_-\) \\
\textbf{WE-2} & Due di tre punti consecutivi in zona gialla o peggio &
Deriva confermata --- intervenire entro un trimestre \\
\textbf{WE-3} & Quattro di cinque punti consecutivi in zona gialla &
Trend strutturale --- riallineare leve operative del Cap. 5.5 \\
\textbf{WE-4} & Otto punti consecutivi in derivata \(\dot d < 0\) &
Erosione lenta --- diagnosi delle relazioni R1-R6 più stressate \\
\textbf{WE-5} & Inversione di derivata seconda \(\ddot d < 0\) con
\(d > 0\) & Pre-allarme --- sistema ancora sicuro ma sta accelerando il
declino \\
\end{longtable}}
}

\Proved{} \textbf{Test sui dati 2024-2026 dei tre attori} (verifica
retrospettiva qualitativa):

{\def\LTcaptype{none} % do not increment counter
{\small\begin{longtable}[]{@{}
  >{\raggedright\arraybackslash}p{(\linewidth - 4\tabcolsep) * \real{0.3333}}
  >{\raggedright\arraybackslash}p{(\linewidth - 4\tabcolsep) * \real{0.3333}}
  >{\raggedright\arraybackslash}p{(\linewidth - 4\tabcolsep) * \real{0.3333}}@{}}
\toprule\noalign{}
\begin{minipage}[b]{\linewidth}\raggedright
Attore
\end{minipage} & \begin{minipage}[b]{\linewidth}\raggedright
Pattern attivati 2024-2026
\end{minipage} & \begin{minipage}[b]{\linewidth}\raggedright
Coerenza con Cap. 7-9
\end{minipage} \\
\midrule\noalign{}
\endhead
\bottomrule\noalign{}
\endlastfoot
USA & WE-1 (zona rossa, V-Dem \(0{,}57\)), WE-3 (Pew trust in declino 4
anni), WE-4 (erosione lenta su tutti i pannelli) & \Proved{} Cap. 7 \\
Cina & WE-4 su P2 (deflazione PPI 10 trim.), WE-2 su P4 (involution +
property) & \Proved{} Cap. 8 \\
Brasile & WE-5 su pannello aggregato (approval calante, ma \(d > 0\)) &
\Proved{} Cap. 9.5 \\
\end{longtable}}
}

\Hypoth{} \textbf{Conseguenza metodologica.} I tre attori esibiscono
pattern WE-DEQ$^2$ \emph{qualitativamente diversi} --- il che valida la
pertinenza diagnostica del cruscotto. La diagnosi non sostituisce
l'interpretazione strutturale dei Cap. 7-9 ma la \emph{quantifica} in
segnali operativi standard.



\section{Scenari Quantizzati per Attore:
2026-2031}\label{scenari-quantizzati-per-attore-2026-2031}

\Hypoth{} \textbf{Discretizzazione per scenari.} Anziché simulare il
sistema dinamico continuo del Cap. 5.4, adottiamo \emph{scenari
quantizzati}: tre traiettorie per attore (best/median/worst case)
parametrizzate dal valore di una variabile critica esogena (per gli USA:
\(\zeta_{\text{Sol}}(t)\); per la Cina: \(K^{\text{imposto}}\); per il
Brasile: esito elettorale 4/10/2026). Le tre traiettorie partono dal
punto stimato \(\boldsymbol{p}(2026)\) del Cap. 5.2 e si sviluppano in
sei tappe annuali.

\subsection{USA --- Scenari 2026-2031}\label{usa-scenari-2026-2031}

{\def\LTcaptype{none} % do not increment counter
{\small\begin{longtable}[]{@{}
  >{\raggedright\arraybackslash}p{(\linewidth - 6\tabcolsep) * \real{0.2500}}
  >{\raggedright\arraybackslash}p{(\linewidth - 6\tabcolsep) * \real{0.2500}}
  >{\raggedright\arraybackslash}p{(\linewidth - 6\tabcolsep) * \real{0.2500}}
  >{\raggedright\arraybackslash}p{(\linewidth - 6\tabcolsep) * \real{0.2500}}@{}}
\toprule\noalign{}
\begin{minipage}[b]{\linewidth}\raggedright
Anno
\end{minipage} & \begin{minipage}[b]{\linewidth}\raggedright
Scenario \emph{median} (probabilità \(\approx 0{,}50\))
\end{minipage} & \begin{minipage}[b]{\linewidth}\raggedright
Scenario \emph{worst} (\(\approx 0{,}30\))
\end{minipage} & \begin{minipage}[b]{\linewidth}\raggedright
Scenario \emph{best} (\(\approx 0{,}20\))
\end{minipage} \\
\midrule\noalign{}
\endhead
\bottomrule\noalign{}
\endlastfoot
2026 & \(d \approx -0{,}10\) (zona rossa) & \(d \approx -0{,}10\) &
\(d \approx -0{,}10\) \\
2027 & \(d \approx -0{,}15\) ; midterm split, \(\Omega \uparrow\) &
\(d \approx -0{,}25\) ; escalation costituzionale &
\(d \approx -0{,}05\) ; SCOTUS riequilibra \\
2028 & \(d \approx -0{,}18\) ; presidenziali, \(\Phi\) scende
ulteriormente & \(d \approx -0{,}35\) ; transizione di fase consolidata
& \(d \approx 0\) ; nuovo equilibrio post-elettorale \\
2029 & \(d \approx -0{,}20\) ; stato stazionario decoerente &
\(d \approx -0{,}40\) ; instabilità acuta & \(d \approx +0{,}05\) ;
recupero parziale \(s_{iii}\) \\
2030 & \(d \approx -0{,}22\) & \(d \approx -0{,}45\) ; crisi monetaria
possibile & \(d \approx +0{,}08\) \\
2031 & \(d \approx -0{,}25\) ; \emph{potenza materiale senza fase} &
\(d \approx -0{,}50\) ; rottura strutturale & \(d \approx +0{,}10\) ;
consolidamento del recupero \\
\end{longtable}}
}

\Conj{} \textbf{Predizione DEQ$^2$ mediana}:
\(T_s^{(2),\text{sys},\text{USA}} < 1\) in tutti gli scenari per
l'intero orizzonte 2026-2031. Recupero solo nello scenario best,
condizionale a SCOTUS attivo + nuova amministrazione \(\neq\) Trump 2.0.

\subsection{Cina --- Scenari 2026-2031}\label{cina-scenari-2026-2031}

{\def\LTcaptype{none} % do not increment counter
{\small\begin{longtable}[]{@{}
  >{\raggedright\arraybackslash}p{(\linewidth - 6\tabcolsep) * \real{0.2500}}
  >{\raggedright\arraybackslash}p{(\linewidth - 6\tabcolsep) * \real{0.2500}}
  >{\raggedright\arraybackslash}p{(\linewidth - 6\tabcolsep) * \real{0.2500}}
  >{\raggedright\arraybackslash}p{(\linewidth - 6\tabcolsep) * \real{0.2500}}@{}}
\toprule\noalign{}
\begin{minipage}[b]{\linewidth}\raggedright
Anno
\end{minipage} & \begin{minipage}[b]{\linewidth}\raggedright
Scenario \emph{median} (\(\approx 0{,}45\))
\end{minipage} & \begin{minipage}[b]{\linewidth}\raggedright
Scenario \emph{worst} (\(\approx 0{,}30\))
\end{minipage} & \begin{minipage}[b]{\linewidth}\raggedright
Scenario \emph{best} (\(\approx 0{,}25\))
\end{minipage} \\
\midrule\noalign{}
\endhead
\bottomrule\noalign{}
\endlastfoot
2026 & \(d^{\text{oss}} \approx +0{,}25\) ;
\(d^{\text{spont}} \approx -0{,}05\) & come median & come median \\
2027 & \(d^{\text{oss}}\) stabile ; \(d^{\text{spont}} \approx -0{,}10\)
; deflazione \(> 12\) trim. & \(d^{\text{oss}}\) scende a \(+0{,}10\) ;
tensione Taiwan acuta & \(d^{\text{oss}} \approx +0{,}28\) ;
anti-involution efficace, riforma proprietà \\
2028 & divergenza oss/spont si amplia & possibile evento \emph{bloqueo}
Taiwan & recupero TFP modesto \\
2029 & \(d^{\text{oss}}\) inizia a calare per \emph{fatica del K
imposto} & crisi politica interna & \(d^{\text{oss}}\) stabile,
\(d^{\text{spont}}\) recupera \\
2030 & \(d^{\text{oss}} \approx +0{,}10\) ;
\(d^{\text{spont}} \approx -0{,}20\) & rottura del cluster sincronizzato
& \(d^{\text{oss}} \approx +0{,}30\) \\
2031 & candidato a \emph{transizione catastrofica} nel triennio
successivo & transizione catastrofica avvenuta & metà-soft landing \\
\end{longtable}}
}

\Conj{} \textbf{Predizione DEQ$^2$ mediana}: divario crescente
\(d^{\text{oss}} - d^{\text{spont}}\), con \emph{fatica del \(K\)
imposto} visibile a partire dal 2029. La transizione di fase cinese ha
probabilità \(\approx 0{,}40\)-\(0{,}55\) di attivarsi nel triennio
2030-2032 --- \emph{fuori} dall'orizzonte centrale del libro ma dentro
alla coda di rilevanza.

\subsection{Brasile --- Scenari 2026-2031 con Biforcazione 4 ottobre
2026}\label{brasile-scenari-2026-2031-con-biforcazione-4-ottobre-2026}

\Hypoth{} \textbf{Schema doppio.} Per il Brasile gli scenari sono
\emph{condizionali} all'esito elettorale. Dividiamo in due rami
principali:

\textbf{Ramo A/B --- Vittoria Lula o successore di sinistra}
(probabilità aprile 2026 \(\approx 0{,}55\)-\(0{,}70\)):

{\def\LTcaptype{none} % do not increment counter
{\small\begin{longtable}[]{@{}
  >{\raggedright\arraybackslash}p{(\linewidth - 6\tabcolsep) * \real{0.2500}}
  >{\raggedright\arraybackslash}p{(\linewidth - 6\tabcolsep) * \real{0.2500}}
  >{\raggedright\arraybackslash}p{(\linewidth - 6\tabcolsep) * \real{0.2500}}
  >{\raggedright\arraybackslash}p{(\linewidth - 6\tabcolsep) * \real{0.2500}}@{}}
\toprule\noalign{}
\begin{minipage}[b]{\linewidth}\raggedright
Anno
\end{minipage} & \begin{minipage}[b]{\linewidth}\raggedright
Scenario \emph{median}
\end{minipage} & \begin{minipage}[b]{\linewidth}\raggedright
\emph{Worst-A/B}
\end{minipage} & \begin{minipage}[b]{\linewidth}\raggedright
\emph{Best-A/B}
\end{minipage} \\
\midrule\noalign{}
\endhead
\bottomrule\noalign{}
\endlastfoot
2027 & \(d \approx +0{,}12\) ; consolidamento istituzionale &
\(d \approx +0{,}05\) ; fiscale teso & \(d \approx +0{,}20\) ; TFFF
parte forte \\
2028 & \(d \approx +0{,}15\) ; BRICS BMG operativo &
\(d \approx +0{,}08\) & \(d \approx +0{,}25\) ; \emph{$\Phi$-engine}
riconosciuto \\
2029 & \(d \approx +0{,}18\) & \(d \approx +0{,}10\) &
\(d \approx +0{,}28\) \\
2030 & \(d \approx +0{,}20\) ; presidenziali fisiologiche &
\(d \approx +0{,}05\) ; alternanza stretta & \(d \approx +0{,}30\) \\
2031 & \(d \approx +0{,}22\) ; \emph{$\Phi$-engine} consolidato &
\(d \approx +0{,}10\) ; condizionale alla nuova amministrazione &
\(d \approx +0{,}32\) ; leadership globale Sud \\
\end{longtable}}
}

\textbf{Ramo C --- Vittoria Flávio Bolsonaro o destra radicale}
(probabilità \(\approx 0{,}30\)-\(0{,}40\)):

{\def\LTcaptype{none} % do not increment counter
{\small\begin{longtable}[]{@{}
  >{\raggedright\arraybackslash}p{(\linewidth - 6\tabcolsep) * \real{0.2500}}
  >{\raggedright\arraybackslash}p{(\linewidth - 6\tabcolsep) * \real{0.2500}}
  >{\raggedright\arraybackslash}p{(\linewidth - 6\tabcolsep) * \real{0.2500}}
  >{\raggedright\arraybackslash}p{(\linewidth - 6\tabcolsep) * \real{0.2500}}@{}}
\toprule\noalign{}
\begin{minipage}[b]{\linewidth}\raggedright
Anno
\end{minipage} & \begin{minipage}[b]{\linewidth}\raggedright
Scenario \emph{median-C}
\end{minipage} & \begin{minipage}[b]{\linewidth}\raggedright
\emph{Worst-C}
\end{minipage} & \begin{minipage}[b]{\linewidth}\raggedright
\emph{Best-C}
\end{minipage} \\
\midrule\noalign{}
\endhead
\bottomrule\noalign{}
\endlastfoot
2027 & \(d \approx +0{,}05\) ; primi 100 giorni transizione &
\(d \approx -0{,}05\) ; smantellamento STF tentato &
\(d \approx +0{,}10\) ; vincoli istituzionali tengono \\
2028 & \(d \approx -0{,}05\) ; uscita Group of Friends; V-Dem cala
\(0{,}05\) & \(d \approx -0{,}15\) ; rottura BRICS &
\(d \approx +0{,}08\) ; resilienza istituzionale \\
2029 & \(d \approx -0{,}10\) ; \emph{solitonizzazione discorsiva} &
\(d \approx -0{,}25\) ; pattern USA-like & \(d \approx +0{,}05\) \\
2030 & \(d \approx -0{,}12\) & \(d \approx -0{,}30\) & \(d \approx 0\) ;
pareggio precario \\
2031 & \(d \approx -0{,}15\) ; \emph{$\Phi$-engine globale assente} &
\(d \approx -0{,}35\) ; transizione di fase regionale &
\(d \approx -0{,}05\) ; mantenimento basale \\
\end{longtable}}
}

\Hypoth{} \textbf{Tesi forte.} Il \textbf{divario fra Ramo A/B e Ramo C}
nel 2031 è \(\Delta d \approx 0{,}30\)-\(0{,}40\) --- il più ampio fra
qualunque variabile considerata in tutto il libro. È la \textbf{misura
quantitativa dell'importanza dell'elezione brasiliana} del 4 ottobre
2026 per l'intero sistema globale.



\section{Lo Scenario Aggregato 2031: Quattro Mondi
Possibili}\label{lo-scenario-aggregato-2031-quattro-mondi-possibili}

\Hypoth{} \textbf{Costruzione.} Combinando gli scenari mediani per
attore con i due rami brasiliani, otteniamo quattro \emph{mondi
possibili} al 2031, ciascuno caratterizzato da una distribuzione
\(\{d_{\text{USA}}, d_{\text{Cina}}, d_{\text{Brasile}}\}\):

{\def\LTcaptype{none} % do not increment counter
{\small\begin{longtable}[]{@{}
  >{\raggedright\arraybackslash}p{(\linewidth - 10\tabcolsep) * \real{0.1667}}
  >{\raggedright\arraybackslash}p{(\linewidth - 10\tabcolsep) * \real{0.1667}}
  >{\raggedright\arraybackslash}p{(\linewidth - 10\tabcolsep) * \real{0.1667}}
  >{\raggedright\arraybackslash}p{(\linewidth - 10\tabcolsep) * \real{0.1667}}
  >{\raggedright\arraybackslash}p{(\linewidth - 10\tabcolsep) * \real{0.1667}}
  >{\raggedright\arraybackslash}p{(\linewidth - 10\tabcolsep) * \real{0.1667}}@{}}
\toprule\noalign{}
\begin{minipage}[b]{\linewidth}\raggedright
Mondo
\end{minipage} & \begin{minipage}[b]{\linewidth}\raggedright
USA
\end{minipage} & \begin{minipage}[b]{\linewidth}\raggedright
Cina
\end{minipage} & \begin{minipage}[b]{\linewidth}\raggedright
Brasile
\end{minipage} & \begin{minipage}[b]{\linewidth}\raggedright
Probabilità soggettiva
\end{minipage} & \begin{minipage}[b]{\linewidth}\raggedright
Caratteristica strutturale
\end{minipage} \\
\midrule\noalign{}
\endhead
\bottomrule\noalign{}
\endlastfoot
\textbf{M1: Decoerenza globale con \(\Phi\)-engine attivo} & \(-0{,}25\)
& \(+0{,}10/-0{,}20\) & \(+0{,}22\) (A/B) &
\(\approx 0{,}35\)-\(0{,}40\) & Brasile come polo di compensazione del
declino USA-Cina \\
\textbf{M2: Decoerenza globale senza \(\Phi\)-engine} & \(-0{,}25\) &
\(+0{,}10/-0{,}20\) & \(-0{,}15\) (C) & \(\approx 0{,}20\)-\(0{,}25\) &
Multipolarità degenere --- frammentazione senza centro coerente \\
\textbf{M3: Rottura cinese precoce + \(\Phi\)-engine attivo} &
\(-0{,}25\) & \(-0{,}30\) (worst) & \(+0{,}22\) (A/B) &
\(\approx 0{,}12\)-\(0{,}18\) & Riassetto sistemico veloce, Brasile
riferimento \\
\textbf{M4: Recupero USA + Cina soft landing + \(\Phi\)-engine attivo} &
\(+0{,}10\) (best) & \(+0{,}30\) (best) & \(+0{,}22\) (A/B) &
\(\approx 0{,}05\)-\(0{,}10\) & Stabilizzazione plurale --- scenario
meno probabile ma operativamente più stabile \\
\end{longtable}}
}

\Hypoth{} \textbf{Lettura.} I mondi M1 e M2 differiscono \emph{solo} per
il ramo brasiliano. La loro probabilità combinata
(\(\approx 0{,}55\)-\(0{,}65\)) misura la \emph{probabilità che il
sistema globale 2031 sia in regime di decoerenza generalizzata}; la
presenza o assenza di un \(\Phi\)-engine ne determina la
\emph{governabilità}. La Parte V del libro (Cap. 12-13) costruisce
protocolli antifragili condizionati alla distinzione M1/M2.

\Proved{} \textbf{Risultato sintetico.}
\(P(\text{sistema globale 2031 in regime di decoerenza}) \approx 0{,}80\)-\(0{,}85\).
\(P(\text{sistema globale 2031 con almeno un $\Phi$-engine}) \approx 0{,}55\)-\(0{,}70\).
La differenza misura l'\textbf{ampiezza dello spazio di manovra residuo
per il Sud globale e per le strategie del Cap. 12}.



\section{\texorpdfstring{Indice \(C_{\text{DEQ}}^2\) e Confronto
Cross-Attore}{Indice C\_\{\textbackslash text\{DEQ\}\}\^{}2 e Confronto Cross-Attore}}\label{indice-c_textdeq2-e-confronto-cross-attore}

\Hypoth{} \textbf{Definizione operativa.}

\[
C_{\text{DEQ}}^2(\boldsymbol{p}, t) := \frac{d(\boldsymbol{p}(t), \mathcal{S})}{\sigma_d(t)}
\]

dove \(\sigma_d(t)\) è la deviazione standard naturale della distanza
calcolata su una finestra mobile di 36 mesi precedenti \(t\).

\Proved{} \textbf{Interpretazione.} \(C_{\text{DEQ}}^2 \geq 1\): il
sistema è a oltre una deviazione standard dalla soglia critica ---
\emph{capability} sistemica accettabile. \(C_{\text{DEQ}}^2 \leq 0\): il
sistema è dentro \(\mathcal{R}_-\). Il valore tipico desiderato per
sistemi industriali è \(C_p \geq 1{,}33\); per sistemi sociali, anche
per la maggiore stochasticity intrinseca, una soglia operativa
ragionevole è \(C_{\text{DEQ}}^2 \geq 1{,}0\).

{\def\LTcaptype{none} % do not increment counter
{\small\begin{longtable}[]{@{}
  >{\raggedright\arraybackslash}p{(\linewidth - 4\tabcolsep) * \real{0.3333}}
  >{\raggedright\arraybackslash}p{(\linewidth - 4\tabcolsep) * \real{0.3333}}
  >{\raggedright\arraybackslash}p{(\linewidth - 4\tabcolsep) * \real{0.3333}}@{}}
\toprule\noalign{}
\begin{minipage}[b]{\linewidth}\raggedright
Attore
\end{minipage} & \begin{minipage}[b]{\linewidth}\raggedright
\(C_{\text{DEQ}}^2\) stimato 2026
\end{minipage} & \begin{minipage}[b]{\linewidth}\raggedright
Lettura
\end{minipage} \\
\midrule\noalign{}
\endhead
\bottomrule\noalign{}
\endlastfoot
USA & \(\approx -0{,}5\) & sotto soglia, dentro \(\mathcal{R}_-\) \\
Cina & \(\approx +1{,}1\) apparente / \(\approx +0{,}3\) effettivo &
apparenza stabile, struttura fragile \\
Brasile & \(\approx +1{,}3\) & unica configurazione
\emph{industrial-grade} \\
\end{longtable}}
}

\Hypoth{} \textbf{Conseguenza.} Il Brasile è l'unico attore della Parte
III a soddisfare il criterio di \emph{capability industriale} applicato
ai sistemi sociali --- coerente con la definizione di \(\Phi\)-engine
del Cap. 9.



\section{Sintesi del Capitolo}\label{sintesi-del-capitolo}

\Proved{} \textbf{In una frase:} la Dashboard SPC-DEQ traduce
\(\boldsymbol{p}(t) \in \mathbb{V}^{(4)}\) in sei pannelli bivariati
monitorabili, un pannello aggregato \(d(\boldsymbol{p}, \mathcal{S})\),
cinque regole di pattern detection adattate da Western Electric, e un
indice scalare \(C_{\text{DEQ}}^2\) --- fornendo l'interfaccia operativa
per le strategie del Cap. 12 e il quadro evidenziale per le predizioni
del Cap. 11. I quattro mondi possibili al 2031 (M1-M4) consolidano
l'analisi della Parte III in \emph{scenari probabilizzati} con
probabilità soggettive esplicite.

{\def\LTcaptype{none} % do not increment counter
{\small\begin{longtable}[]{@{}
  >{\raggedright\arraybackslash}p{(\linewidth - 4\tabcolsep) * \real{0.3333}}
  >{\raggedright\arraybackslash}p{(\linewidth - 4\tabcolsep) * \real{0.3333}}
  >{\raggedright\arraybackslash}p{(\linewidth - 4\tabcolsep) * \real{0.3333}}@{}}
\toprule\noalign{}
\begin{minipage}[b]{\linewidth}\raggedright
Strumento
\end{minipage} & \begin{minipage}[b]{\linewidth}\raggedright
Funzione
\end{minipage} & \begin{minipage}[b]{\linewidth}\raggedright
Capitolo che lo usa
\end{minipage} \\
\midrule\noalign{}
\endhead
\bottomrule\noalign{}
\endlastfoot
Sei pannelli bivariati P1-P6 & Diagnosi delle relazioni R1-R6 più
stressate & Cap. 12 (interventi mirati) \\
Pannello aggregato \(P_0 = d(\boldsymbol{p}, \mathcal{S})\) & Indicatore
scalare di stabilità & Cap. 11 (predizioni) \\
Cinque regole WE-DEQ$^2$ & Pattern detection precoce & Cap. 12 (allarme
operativo) \\
Indice \(C_{\text{DEQ}}^2\) & Confronto cross-attore & Cap. 11 (sintesi
quantitativa) \\
Scenari quantizzati per attore & Proiezioni 2026-2031 & Cap. 11
(predizioni datate) \\
Quattro mondi possibili 2031 & Scenari aggregati & Cap. 12-13 (strategie
condizionali) \\
\end{longtable}}
}

\Hypoth{} \textbf{Tesi conclusiva.} La DEQ$^2$ è ora \emph{operativa}. La
Parte IV chiude la fase teorica del libro e apre la fase normativa: dato
lo spazio degli scenari, \emph{quali strategie sono antifragili in
ciascuno di essi?} Questo è l'oggetto della Parte V.



\section{Note Epistemiche di Chiusura
Capitolo}\label{note-epistemiche-di-chiusura-capitolo}

{\def\LTcaptype{none} % do not increment counter
{\small\begin{longtable}[]{@{}
  >{\raggedright\arraybackslash}p{(\linewidth - 4\tabcolsep) * \real{0.3333}}
  >{\raggedright\arraybackslash}p{(\linewidth - 4\tabcolsep) * \real{0.3333}}
  >{\raggedright\arraybackslash}p{(\linewidth - 4\tabcolsep) * \real{0.3333}}@{}}
\toprule\noalign{}
\begin{minipage}[b]{\linewidth}\raggedright
\#
\end{minipage} & \begin{minipage}[b]{\linewidth}\raggedright
Affermazione
\end{minipage} & \begin{minipage}[b]{\linewidth}\raggedright
Status
\end{minipage} \\
\midrule\noalign{}
\endhead
\bottomrule\noalign{}
\endlastfoot
1 & Adattamento SPC industriale $\to$ DEQ$^2$ & \Hypoth{} traduzione
strutturale rigorosa \\
2 & Sei pannelli bivariati con iso-curve esplicite & \Hypoth{}
costruzione operativa \\
3 & Cinque regole WE-DEQ$^2$ & \Hypoth{} adattamento di Western Electric
(1956) \\
4 & Test retrospettivo 2024-2026 sui tre attori & \Proved{} pattern
coerenti con Cap. 7-9 \\
5 & Scenari quantizzati USA con \(d \to -0{,}25\) entro 2031 (median) &
\Conj{} congettura datata, derivata da Cap. 7 \\
6 & Scenari Cina con divergenza \(d^{\text{oss}}/d^{\text{spont}}\)
crescente & \Conj{} congettura, derivata da Cap. 8 \\
7 & Doppio ramo Brasile (A/B vs C) con
\(\Delta d \approx 0{,}30\)-\(0{,}40\) al 2031 & \Conj{} quantificazione
della biforcazione del Cap. 9.5 \\
8 & Quattro mondi possibili al 2031 con probabilità soggettive & \Conj{}
stime soggettive su evidenza Cap. 7-9 \\
9 & \(P(\text{decoerenza globale 2031}) \approx 0{,}80\)-\(0{,}85\) &
\Conj{} sintesi probabilistica \\
10 & Indice \(C_{\text{DEQ}}^2\) + soglia operativa \(\geq 1{,}0\) &
\Hypoth{} adattamento di \(C_p\) industriale \\
11 & Solo Brasile soddisfa \(C_{\text{DEQ}}^2 \geq 1{,}0\) & \Hypoth{}
conseguenza derivata coerente con Cap. 9 \\
\end{longtable}}
}


